\begin{ledgroupsized}[r]{120mm}
\footnotesize
\pstart
\noindent\textbf{Überlieferung:}
\pend
\end{ledgroupsized}
\begin{ledgroupsized}[r]{114mm}
\footnotesize
\pstart
\parindent -6mm
\makebox[6mm][l]{\textit{L}}%
Konzept:
LH~XXXVII~4~Bl.~2. 
Ein Blatt~4\textsuperscript{o} (ca.~19~x~20~cm.);		%
\href{\Scan{LH_37_04_002wz}}%
{Wasserzeichenfragment} am Blattrand;
Papiererhaltungsmaßnahmen.
Zwei Seiten.
\pend
\end{ledgroupsized}
\selectlanguage{latin}		%
\vspace{8mm}
\pstart
\noindent
\href{\Scan{LH_37_04_002r}}%
{\lbrack2~r\textsuperscript{o}\rbrack} 
\pend
\nobreak 
\begin{edfigure}[-.5\baselineskip][0em]
\includegraphics[scale=1]{%
gesamttex/edit_VIII-4/images/LH_37_04_002_d1_002r.pdf}
\edcaption[0em]{%
\edtext{\lbrack\textit{Fig.~1}\rbrack}{%
\lemma{\bracks{\textit{Fig.~1}}}%
\Cfootnote[nonum]{%
Zwei gestrichene und zwei nicht gestrichene Entwürfe zum Diagramm werden nicht wiedergegeben.}%
\lemma{}%
\Afootnote[nonum]{%
\!\textit{Oberhalb \bracks{\textit{Fig.~1}}}: In casu simplicis descensus per planum inclinatum\textsuperscript{\bracks{a}} ex theoremate nostro novo nascetur theorema vulgare.
\textsuperscript{\bracks{b}}
\\[1.5ex]
{\footnotesize 
\textsuperscript{\bracks{a}}inclinatum %
\textit{(1)}~pro %
\textit{(2)}~ex~\textit{L}
\qquad
\textsuperscript{\bracks{b}}theorema vulgare: Siehe bspw.
\indexN{\textls{Galilei} (Galilaeus, Galileus), Galileo 1564\textendash1642}%
\textsc{G.~Galilei}, %
\cite{00050}\textit{Discorsi}, %
Giornata Terza; \glqq Aggiunta postuma dell'Autore\grqq\ %
nach dem Scholion zum 2.~Theorem. Der Zusatz erschien erst in der posthumen Ausgabe: \cite{01084}\textit{Opere}, Bologna 1656, Bd.~II, S.\;132\lbrack ter\rbrack\textendash134 (\cite{00048}\textit{GO} VIII, S.\;216\textendash219 im Apparat).}}%
\lemma{}%
\Afootnote[nonum]{%
\!\textit{Unterhalb \bracks{\textit{Fig.~1}}, gestrichen}:
{\footnotesize Poterunt puncta \textit{A} et \textit{B} sumi coincidentia,\textsuperscript{\bracks{a}} seu in intersectione plani dati \textit{AD}, et rectae positione datae}
\\[1.5ex]
{\footnotesize 
\textsuperscript{\bracks{a}}coincidentia, \textit{(1)}~ad \textit{(2)}~seu~\textit{L}}}}%
}
\end{edfigure}
\clearpage
\pstart
\noindent
Sit 
\edtext{globus\indexS{globus} cujus centrum figurae\indexS{centrum figurae} \textit{C}}{%
\lemma{globus}%
\Bfootnote{%
\textit{(1)}~\textit{C} %
\textit{(2)}~eodem sem %
\textit{(3)}~cujus centrum figurae \textit{C}~\textit{L}}}
descendens\indexS{globus descendens in plano sine volutione} in plano\indexS{planum inclinatum} \textit{AD} sine volutione\indexS{volutio} ita ut eodem\indexS{globus eodem puncto planum tangens} semper puncto \textit{D} planum tangat. Ex centro ejus \textit{C} prodeat chorda \textit{CB}, ad trochleam\indexS{trochlea} \textit{B}, per quam globum\indexS{pondus sustinens globum per trochleam} sustinere conetur pondus \textit{E}. Sit 
\edtext{globi centrum gravitatis\indexS{centrum gravitatis a centro figurae diversum} \textit{G}}{%
\lemma{globi}%
\Bfootnote{%
\textit{(1)}~centrum gravitatis %
\textit{(2)}~centrum Gravitatis %
\textit{(a)}~\textit{C} %
\textit{(b)}~\textit{G}~\textit{L}}}
a centro figurae\indexS{centrum figurae} diversum situm in recta \textit{CD} et centrum commune\indexS{centrum commune duorum ponderum}
tam ponderis \textit{E}, quam centri grav.\indexS{centrum gravitatis} \textit{G}, sit \textit{K}\bracks{;}
ex \textit{D}, \textit{G}, \textit{C}, \textit{K} in rectam \textit{BE} demittantur normales \textit{DM}, \textit{GL},
\textit{CH}, \textit{KP}; quaeritur situs globi et ponderis talis, ut \textit{BP} sit omnium possibilium 
minima.
\pend
\pstart
Itaque \textit{BP} per calculum investigetur. Dantur puncta \textit{A} et \textit{B} adeoque recta \textit{AB}. Datur et ratio laterum trianguli \textit{AMD} 
\edtext{vel \textit{DM$\beta$}. Dantur et rectae \textit{CG}, \textit{GD}, \textit{DC}. Datur et summa}{%
\lemma{vel \textit{DM$\beta$}.}%
\Bfootnote{%
\textit{(1)}~Datur et recta \textit{CG} %
\textit{(2)}~Dantur et rectae \textit{CG}, \textit{GD}, \textit{DC}. %
\textit{(a)}~Ut igitur investigemus rectam \textit{BP}, adeoque punctum %
\textit{(b)}~Daturque %
\textit{(c)}~Datur %
\textit{(d)}~Datur \bars{et \textit{erg.}} summa~\textit{L}}}
rectarum \textit{CB}, et \textit{BE}\bracks{,} denique datur
\edtext{ratio}{%
\lemma{}%
\Bfootnote{%
ratio \textit{erg.~L}}}
ponderis globi ad pondus \textit{E}. Assumatur 
\edtext{jam \textit{BM} vel \textit{AM} tanquam esset data,}{%
\lemma{jam}%
\Bfootnote{%
\textit{(1)}~\textit{BH} %
\textit{(2)}~\textit{BM} %
\textit{(a)}~tanquam %
\textit{(b)}~vel \textit{AM} tanquam %
\textit{(aa)}~data %
\textit{(bb)}~esset data,~\textit{L}}}
hinc daretur \textit{DM}, \textit{AC}, \textit{M$\beta$}, \textit{D$\beta$}. Adeoque ob datas \textit{DG}, \textit{DC}, darentur \textit{BG}, 
\edtext{\textit{BC}, adeoque}{%
\lemma{\textit{BC},}%
\Bfootnote{%
\textit{(1)}~eo
\textit{(2)}~adeoque~\textit{L}}}
et \textit{GL}, \textit{CH}, et \textit{$\beta$H}, \textit{$\beta$L} adeoque \textit{BL}, \textit{BH}. Datur igitur et \textit{CB}, adeoque et $CB+BH$, unde si auferas $CB+BE$ dat. restabit \textit{BE}, datur ergo et \textit{EL}; ejusque certa portio \textit{EP}, quae addita ipsa ipsi \textit{BE}, dat \textit{BP}, quae facienda omnium possibilium 
\edtext{maxima.
Haec jam calculo}{%
\lemma{maxima.}%
\Bfootnote{%
\textit{(1)}~Hic calcul %
\textit{(2)}~Haec jam %
\textit{(a)}~data %
\textit{(b)}~calculo~\textit{L}}}
exequamur. %
\pend
\pstart
\textit{AB}, \textit{a}. \textit{BR}, \textit{b}. \textit{AR}, 
\edtext{$c=\sqrt{aa+bb}$.}{\lemma{}\Bfootnote{$=\sqrt{aa+bb}$ \textit{erg.~L}}}
\textit{CD}, \textit{d}. \textit{CG}, \textit{g}. \textit{GD}, $d-g$. 
\edtext{$CB+BE=e$. Globus}{%
\lemma{$CB+BE=e$.}%
\Bfootnote{%
\textit{(1)}~\textit{C} ad
\textit{(2)}~Globus~\textit{L}}}
\textit{C} ad pondus \textit{E} sit \textit{m} ad \textit{n}.
\pend
\pstart
\textit{AM}, \textit{y}. \textit{DM}\bracks{,} $yb:a$.
\textit{AD}, $yc:a$.
\textit{M$\beta$}\bracks{,} $ybb:aa$. \textit{D$\beta$}, $ybc:aa$. Et $\beta G=\overline{ybc:aa}+d-g$ et $GL,=\overline{\overline{ybc:aa}+d-g}\; a:c$ et $L\beta =\overline{\overline{ybc:aa}+d-g}\; b:c$ et $\beta C=ybc:aa+d$ et $CH=\overline{\overline{ybc:aa}+d}\; a:c$ 
et $H\beta =\overline{\overline{ybc:aa}+d}\;$
\edtext{$b:c$ et 
\edtext{$AH,y-\overline{db:c}$}{%
\lemma{}%
\Afootnote{%
\textit{Am Rand}: $HM:CD\squaredots b:c$ seu \textit{HM}, $db:c$}}
et $BH,y-\overline{db:c}-a$.}{%
\lemma{$b:c$}%
\Bfootnote{%
\textit{(1)}~$BH=y-a$ %
\textit{(2)}~$CB=\sqrt{yy}$ %
\textit{(3)}~et $AH,y-\overline{db:c}$ et $BH,y-\overline{db:c}-a$.~\textit{L}}}
\pend
\pstart
$CB^2=CH^2+BH^2$, jam $CH\!=$
\edtext{$yb:a+da:c$ et \textit{EH}}{%
\lemma{$yb:a+da:c$}%
\Bfootnote{%
\textit{(1)}~et $CH^2=yybb:aa+2ybd:c+ddaa:cc$ %
\textit{(2)}~\bars{et \textit{streicht Hrsg.}}  $BH^2=yy-2ya+aa$. 
\bars{Fiet $\overline{B}^2\!=$ \textit{streicht Hrsg.}} $yy-\overline{\overline{bb:aa}+1}+2y\overline{\overline{bd:c}-a}+aa\,\overline{\overline{dd:cc}-1}$. %
\textit{(a)}~Unde auferen %
\textit{(b)}~A \textit{CB}, auferendo $CB+BE=e$, fiet %
\textit{(c)}~et $CB+BH-$ %
\textit{(aa)}~$\underbrace{\dfrac{CB-BE}{}}_{e}$ %
\textit{(bb)}~$\overbrace{CB-BE}^{e}=\sqrt{yy-\overline{\overline{bb:aa}+1}+2y\;\overline{\overline{bd:c}-a}+aa\,\overline{\overline{dd:cc}-1}}+y-a-e=EH$. Porro %
\textit{(cc)}~et \textit{EH}~\textit{L}}}
\edtext{%
$=CB+BH\overbrace{-CB-BE}^{-e}=\sqrt{\overline{\overline{yb:a}+\overline{da:c}}^2+\overline{y-\overline{db:c}-a}^2}+y-\overline{db:c}-a-e$.%
}{%
\lemma{}%
\Afootnote{%
\textit{Am Rand}: \textsuperscript{\lbrack a\rbrack} $d\overline{x^e}=ex^{\underline{e-1}}$\bracks{,} $dx^{1:2}=\dfrac{1}{2}dx\; x^{-1:2}=\dfrac{dx}{2\surd x}$\bracks{,} $dv^2=2vdv$
\\[1.5ex]
{\footnotesize \textsuperscript{\lbrack a\rbrack}%
\textit{(1)}~$dx^2=2xdx$. $d\,\overline{\surd x}=\dfrac{dx}{2\surd x}$. %
\textit{(2)}~\bars{$dx^e=$ \textit{streicht Hrsg.}} $d\overline{x}\; x^{\underline{e-1}}:e$. %
\textit{(3)}~$d\overline{x^e}=ex^{\underline{e-1}}$~\textit{L}}}}
$HL:HM\squaredots CG$ seu $g:CD$ seu \textit{d} ergo $HL=\overline{db:c}\; g:d=bg:c$ et $EL=EH+HL,=EH+\overline{bg:c}$. Jam $EP:EL\squaredots m:\overline{m+n}$. Fit $EP=\overline{EH+\overline{bg:c}}\; m:\overline{m+n}$ et $BE=e-CB$ et $BP=BE+EP=e-CB+\overline{EH+\overline{bg:c}}\; m:\overline{m+n}=e-CB+\overline{CB-e+BH+\overline{bg:c}}\; m:\overline{m+n}$. Ergo $BP=\overline{e-CB}\; m:\overline{m+n}+BH\; m:\overline{m+n}+bgm:\overline{c\cdot\overline{m+n}}=$ Maxima.
\pend
\pstart
Jam $d\,\overline{CB}=\overline{\overline{yb:a}+\overline{da:c}}\; d\,\overline{y}\,{b:a}\;+\overline{y-db:c-a}$
\edlabel{37_04_002_1a}%
\edtext{}{%
\xxref{37_04_002_1a}{37_04_002_1b}%
\lemma{$dy\; y:CB$.}%
\Bfootnote{%
\textit{(1)}~Sed %
\textit{(2)}~Seu $=$ %
\textit{(3)}~Seu $d\,\overline{CB} =$~\textit{L}}}%
$dy\; y:CB$.%
\pend
\pstart
Seu $d\,\overline{CB} =$%
\edlabel{37_04_002_1b}
$\overline{ybb:aa}\;\ovalbox{+\overline{db:c}}+y\;\ovalbox{-db:c}-a,\,\cdot\, d\overline{y}:CB$ et
\edlabel{37_04_002_2a}%
\edtext{}{%
\xxref{37_04_002_2a}{37_04_002_2b}%
\lemma{$d\overline{BH}=d\overline{y}$.}%
\Bfootnote{%
\textit{(1)}~Ergo fit %
\textit{(2)}~Jam~\textit{L}}}%
$d\overline{BH}=d\overline{y}$.
\pend
\pstart
Jam\edlabel{37_04_002_2b}
$d\,\overline{BP}=-d\,\overline{CB}\; n:\overline{m+n}+d\,\overline{BH}\; m:\overline{m+n}=0$. %
\pend
\pstart
Ergo fiet denique $nybb:aa+ny-aa=m\;\overline{CB}$, seu $n:m\squaredots CB:\left(\overbrace{ybb:aa}^{M\beta}+\stackrel{\displaystyle{AM}}{y}-\stackrel{\displaystyle{AB}}{a}=\right)B\beta$. %
\pend
\pstart
\edlabel{37_04_002_3a}%
Ergo habetur \textls{theorema memorabile},\indexS{theorema memorabile} \textit{n} ad \textit{m}, seu pondus \textit{E} ad globum \textit{C}, ut \textit{CB} ad \textit{B$\beta$}.%
\edlabel{37_04_002_3b}
Et valde memorabile est, quod nihil refert ubi sit centrum\indexS{centrum gravitatis} 
\edtext{gravitatis globi}{%
\lemma{gravitatis}%
\Bfootnote{%
\textit{(1)}~modo cadat in rectam \textit{CD} %
\textit{(2)}~cor
\textit{(3)}~globi~\textit{L}}}
seu corporis, modo cadat in rectam \textit{CD}\bracks{,} pondus \textit{E} ad globum \textit{C} ut \textit{CT} ad \textit{TD}. 
Patet etiam nil referre, utrum \textit{CD} sit globus et \textit{C} ejus centrum,\indexS{centrum globi}
tantum enim ex centro gravitatis\indexS{centrum gravitatis} \textit{G} ducatur in planum sustentans\indexS{planum sustentans} normalis \textit{GD} quae producta rectam tendentiae reliquae 
\edtext{sustentationis\indexS{tendentia sustentationis plani} \textit{BC}}{%
\lemma{sustentationis}%
\Bfootnote{%
\textit{(1)}~secet \textit{C} %
\textit{(2)}~\textit{BC}~\textit{L}}}
secet in \textit{C}. Sed hic supposuimus tendentiam sustentationis\indexS{tendentia sustentationis plani} plani \textit{AD} perpendicularem ad planum ex contactu \textit{D} cadere in \textit{G} centrum gravitatis.\indexS{centrum gravitatis}
\pend
\pstart
Quaerendum superest qua vi\indexS{vi qua planum premitur} prematur planum \textit{AD}. Videndum et quid futurum fuisset, si centrum gravitatis\indexS{centrum gravitatis} \textit{G} non cecidisset in rectam \textit{CD}. Nostra demonstratio supponit corpus\indexS{corpus volvi non potens} non posse 
\edtext{volvi ut si affixa ei esset recta \textit{SV}, ubi si \textit{V} esset non perpendiculariter sub \textit{G} sed magis ad dextram versus \textit{D}, posset \textit{G} volvi\indexS{corpus volvi seu rotari potens} seu rotari circa \textit{V}. 
Sed si contra esset \textit{V} praecise sub \textit{G} vel magis ad sinistram ipsius \textit{G} conatus volvendi\indexS{conatus volvendi} impedietur, et perinde erit ac si corpus cujus centrum grav.\indexS{centrum gravitatis} \textit{G} affixum in puncto \textit{C} inniteretur plano recta \textit{XV}\bracks{,} est datum punctum \textit{X} illud quo \textit{GL} horizontalis secat \textit{AD}, nam quod supra est \textit{XS} non sustinet globum.}{%
\lemma{}%
\Bfootnote{%
\textit{(1)}~ut si affixa ei esset recta \textit{SV} %
\textit{(2)}~ut si \bracks{...} \textit{V} esset 
\textit{(a)}~perpendiculariter sub \textit{G} vel
\textit{(b)}~non perpendiculariter sub \textit{G} sed
\bracks{...} globum \textit{erg.}. Sed~\textit{L}}}
Sed operae pretium erit inquisitionem instituere generaliorem.
\pend
\begin{edfigure}%
\includegraphics[scale=1]{%
gesamttex/edit_VIII-4/images/LH_37_04_002_d2_002v.pdf}
\edcaption{\edtext{\lbrack\textit{Fig.~2}\rbrack}{%
\lemma{\lbrack\textit{Fig.~2}\rbrack}%
\Cfootnote[nonum]{%
Ein gestrichener Entwurf zum Diagramm wird nicht wiedergegeben.}}}
\end{edfigure}
\pstart
\href{\Scan{LH_37_04_002v}}%
{\lbrack2~v\textsuperscript{o}\rbrack} 
Consideratu dignum 
\edtext{vide\bracks{tur}}{\lemma{}\Bfootnote{videri \textit{L ändert Hrsg.}}}
omnem sustinendi modum\indexS{modus sustinendi tendens in centrum gravitatis} sic accipiendum ut tendentem in centrum gravitatis.\indexS{centrum gravitatis} Sit enim globus \textit{C}, cujus centrum gravitatis\indexS{centrum gravitatis} \textit{G}, innitens plano inclinato\indexS{planum inclinatum} \textit{AD}, sustentatus filo\indexS{globus filo sustentatus} \textit{HF} quod productum tendat in centrum gravitatis.\indexS{centrum gravitatis}
Hoc filum concipiatur replicatum circa trochleam\indexS{trochlea} \textit{H} atque inde tendere versus aliam trochleam\indexS{trochlea} \textit{K}, atque trahi a suspenso inde pondere\indexS{pondus suspensum} \textit{E}. 
\edtext{Patet nihil referre, quanta sit longitudo}{%
\lemma{}%
\Afootnote{%
\textit{Am Rand}: Potest etiam concipi\textsuperscript{\bracks{a}} ipsa \textit{DL} secante lineam directionis ponderis \textit{G}, in \textit{R}, esse \textit{LR} chordam quae in \textit{R} circa trochleam\indexS{trochlea} jacta trahat \textit{G} perpendiculariter,\indexS{chorda trahens globum perpendiculariter circa trochlea} et elegans erit si quacunque facta hypothesi omnia redirent eodem ut arbitror.
\\[1.5ex] %
{\footnotesize 
\textsuperscript{\bracks{a}}concipi \bars{et \textit{gestr.}} ipsa~\textit{L}}}}
fili \textit{HF}, hac igitur evanescente perinde erit 
\edtext{ac si globus\indexS{globus filo sustentatus} sustineretur}{%
\lemma{ac si}%
\Bfootnote{%
\textit{(1)}~pondus %
\textit{(2)}~globus %
\textit{(a)}~sustinet %
\textit{(b)}~sustineretur~\textit{L}}}
filo \textit{FK}. Ergo vicissim globo sustentato\indexS{globus filo sustentatus} a filo \textit{FK} perinde est ac si sustineretur\indexS{globus filo sustentatus} a filo \textit{HF}\bracks{,} 
adeoque sustentatio\indexS{sustentatio tendens in centrum gravitatis} sumi potest ut tendens in centrum gravitatis\indexS{centrum gravitatis}\indexS{modus sustinendi tendens in centrum gravitatis} 
quod valde videtur memorabile, et ob conclusionem et ob modum argumentandi. 
Sic eodem modo pro plano sustentante\indexS{planum sustentans} \textit{AD}, finge filum \textit{LMND} replicatum transiens per trochleas\indexS{trochlea} \textit{M} et \textit{N} et ex 
\edtext{\bracks{\textit{N}}}{\lemma{}\Bfootnote{\textit{M} \textit{L ändert Hrsg.}}}
tendens in \textit{D}, sic ut \textit{N} cadat in 
\edtext{rectam \textit{GD} ubi sustentatio\indexS{sustentatio tendens in centrum gravitatis}}{%
\lemma{rectam}%
\Bfootnote{%
\bars{inter \textit{gestr.}}
\textit{(2)}~\textit{GD} %
\textit{(1)}~.~Patet si puncta \textit{M}, \textit{N}, \textit{D} %
\textit{(2)}~.~Pona %
\textit{(3)}~ubi %
\textit{(a)}~tendentia %
\textit{(b)}~sustentatio~\textit{L}}}
tendit in centrum gravitatis\indexS{centrum gravitatis} et tamen hic casus coincidit cum sustentatione plani\indexS{sustentatio plani} \textit{AD}, seu cum sustentatione chordae\indexS{sustentatio chordae} \textit{LD} si puncta \textit{D}, \textit{M}, \textit{N} fingantur coincidere seu esse distantiae infinite parvae. Caeterum ut chorda \textit{LD} a corpore
\edtext{\bracks{non}}{\lemma{}\Bfootnote{non \textit{erg.~Hrsg.}}}
impediatur fingi potest infinitae tenuitatis, et corpus perforatum,\indexS{corpus perforatum} 
pro ipsius transitu, detrimento perforationis infinite parvo\indexS{detrimentum perforationis infinite parvum}
seu nil nocente.
Quid si in pluribus punctis\indexS{planum in pluribus punctis sustentans}
\edtext{simul sustineat}{%
\lemma{simul}%
\Bfootnote{%
\textit{(1)}~sustineatur %
\textit{(2)}~sustineat~\textit{L}}}
planum? Poterimus 
\edtext{considerare infinitas}{\lemma{}\Bfootnote{considerare \bars{ut \textit{gestr.}} infinitas~\textit{L}}}
rectas in centrum gravitatis\indexS{centrum gravitatis} tendentes,\indexS{rectae infinitae in centrum gravitatis tendentes} et componendo inde facere unum nisum.\indexS{nisus compositus} 
Sed haec opinor in corporibus revera sustentatis,\indexS{corpus sustentatum} 
nam si partem libertatis habeant, videndum an aliud saltem in ictibus\indexS{ictus} ubi 
eccentricae impressiones\indexS{impressio eccentrica} contribuunt ad gyrationes,\indexS{gyratio} 
an forte hinc inde ipsa gyrationum origo\indexS{origo gyrationum} manifestari potest cum non est aequilibrium.
\pend
\begin{edfigure}%
\includegraphics[scale=1]{%
gesamttex/edit_VIII-4/images/LH_37_04_002_d3_002v.pdf}
\edcaption{\lbrack\textit{Fig.~3}\rbrack}
\end{edfigure}
\pstart
Perfecte consentit 
\edtext{theorema pagina praecedente a me inventum}{%
\lemma{theorema pagina praecedente a me inventum}%
\Cfootnote{Siehe \refpassage{37_04_002_3a}{37_04_002_3b}.}}
\edlabel{37_04_002_4a}%
\edtext{}{%
\xxref{37_04_002_4a}{37_04_002_4b}%
\lemma{cum theoremate \bracks{...} ut \textit{CG}}%
\Cfootnote{Mit "theorema generalissimum" ist die Parallelogramm-Regel gemeint:
vgl. 
\indexN{\textls{Varignon}, Pierre de 1654\textendash1722}\textsc{P.~Varignon},
\cite{02028}\textit{Projet d'une nouvelle mechanique}, Paris 1687, bes. "Proposition fondamentale des poids suspendus avec des cordes", S.\;11; Leibniz hat zwischen Dezember 1689 und 2. Februar 1690 Auszüge angefertigt, siehe N.\;\ref{60365} in diesem Band.}}%
\edtext{cum theoremate}{%
\lemma{cum}%
\Bfootnote{%
\textit{(1)}~alio
\textit{(2)}~theoremate~\textit{L}}}
illo generalissimo
nempe tres sunt directiones\indexS{directiones tres se intersecantes} se intersecantes\bracks{:} potentiae\indexS{directiones ponderis et potentiarum} seu ponderis \textit{G}, nempe \textit{FG}, 
\edtext{potentiae \textit{T} (seu}{%
\lemma{potentiae}%
\Bfootnote{%
\textit{(1)}~seu %
\textit{(2)}~\textit{T} %
\textit{(a)}~ponderis %
\textit{(b)}~(seu~\textit{L}}}
ponderis \textit{E}) nempe \textit{FC} et potentiae resistentis\indexS{potentia resistens} \textit{D}, 
nempe \textit{GC} et erit
\edtext{pondus \textit{E}, seu potentia \textit{T} ad pondus\indexS{pondus} seu potentiam\indexS{potentia} \textit{G}}{%
\lemma{pondus \textit{E},}%
\Bfootnote{%
\textit{(1)}~ad %
\textit{(2)}~ad potentiam \textit{G} seu po %
\textit{(3)}~seu potentia \textit{T} %
\textit{(a)}~ad %
\textit{(b)}~ad~\textit{L}}}
pondus seu potentiam \textit{G} ut \textit{CF} 
\edtext{ad \textit{FG}, pressio\indexS{pressio globi in planum} autem seu resistentia\indexS{resistentia plani}}{%
\lemma{ad \textit{FG},}%
\Bfootnote{%
\textit{(1)}~resistentia autem %
\textit{(2)}~pressio autem seu resistentia~\textit{L}}}
ipsius \textit{D} est ut \textit{CG}.%
\edlabel{37_04_002_4b}
Haec proportio \textit{CF} ad \textit{FG} coincidit cum supra dicta \textit{CT} ad \textit{TD}.
\pend
